\documentclass[11pt]{article}
\usepackage[a4paper,left=0.9in,right=0.9in,top=0.9in,bottom=0.9in]{geometry}
\usepackage[english]{babel}
\usepackage{amsmath,amssymb,mathtools}
\usepackage{algorithm}
\usepackage{algpseudocode}
\usepackage{booktabs}
\usepackage{enumitem}
\usepackage{hyperref}
\usepackage{microtype}
\usepackage{xcolor}

\setlength{\parindent}{1.5em}
\setlist[itemize]{itemsep=0.2em,topsep=0.35em}
\hypersetup{colorlinks=true,linkcolor=blue!55!black,urlcolor=blue!55!black}
\newcommand{\R}{\mathbb{R}}
\newcommand{\conv}{\operatorname{conv}}
\newcommand{\OPT}{\operatorname{OPT}}

\title{\textbf{Branch and Bound for the Fixed-Order\\Nonconvex Touring Polygons Problem}}
\author{Gabriel F. Ushijima}
\date{September 4, 2026}

\begin{document}
\maketitle

\begin{abstract}
This note presents the current exact algorithm for the fixed-order nonconvex
Touring Polygons Problem (TPP). The central idea is to decompose every
nonconvex polygon into convex pieces, branch over the choice of one piece per
polygon, and use the exact convex TPP algorithm to bound partial choices. We
first explain the method through a small example, then give pseudocode and a
proof of correctness, and finally report experiments on the 558 instances
distributed by Fekete, Kniep, Krupke, and Perk (TU Braunschweig). The
experiments identify repeated convex solves as the dominant cost and show that
finding a good feasible path is usually easier than proving its optimality.
\end{abstract}

\section{Problem and main idea}

The input consists of two points $s,t\in\R^2$ and an ordered sequence of
polygons $P_1,P_2,\ldots,P_k$. We seek a shortest path from $s$ to $t$ that
intersects the polygons in this order. Equivalently, we choose contact points
$x_i\in P_i$ and minimize
\[
	\lVert s-x_1\rVert+\sum_{i=1}^{k-1}\lVert x_i-x_{i+1}\rVert+\lVert x_k-t\rVert.
\]
The order is fixed. Thus, this algorithm does not search over the $k!$
permutations of the polygons.

Suppose polygon $P_i$ is decomposed into convex pieces
\[
	P_i=C_{i,1}\cup C_{i,2}\cup\cdots\cup C_{i,m_i}.
\]
A path that visits $P_i$ contains a point in at least one of these pieces.
Consequently, the nonconvex problem can be solved by considering every vector
of piece choices $(j_1,\ldots,j_k)$, where
$j_i\in\{1,\ldots,m_i\}$, and solving the convex TPP instance
$(C_{1,j_1},\ldots,C_{k,j_k})$. Direct enumeration requires $\prod_i m_i$
convex solves, so it is generally infeasible. Branch and bound avoids most of
these complete combinations.

\subsection{A three-polygon example}

Assume the decompositions have respectively $2$, $3$, and $2$ pieces. Direct
enumeration has $2\cdot3\cdot2=12$ leaves. At node $(2,1)$, the first two
choices are fixed and the third polygon is undecided. We solve the relaxed
convex instance
\[
	(C_{1,2},C_{2,1},H_3),\qquad H_3=\conv(P_3).
\]
If this relaxed path has length $14$, while a feasible path of length $13$ is
known, neither $(2,1,1)$ nor $(2,1,2)$ can improve the incumbent. Both leaves
are discarded after one bound computation.

This example contains the complete logic of the method:
\begin{enumerate}
	\item a complete piece choice gives a feasible solution and an upper bound;
	\item replacing undecided polygons by larger convex regions gives a lower bound;
	\item a subtree is discarded when its lower bound is worse than the best
	known feasible solution.
\end{enumerate}

\section{Search tree and bounds}

A node $u=(j_1,\ldots,j_d)$ at depth $d$ fixes one piece for each of the first
$d$ polygons. Its children append one of the $m_{d+1}$ pieces of the next
polygon. The root is the empty vector and a node at depth $k$ is a leaf.

For every polygon, preprocessing constructs
\[
	H_i=\conv\left(\bigcup_{q=1}^{m_i}C_{i,q}\right).
\]
For a partial choice $u=(j_1,\ldots,j_d)$, define
\[
	R(u)=(C_{1,j_1},\ldots,C_{d,j_d},H_{d+1},\ldots,H_k),
	\qquad LB(u)=\OPT(R(u)).
\]
All regions in $R(u)$ are convex, so the existing convex TPP algorithm solves
this instance exactly. The bound is valid because $C_{i,q}\subseteq H_i$.
Replacing a feasible region by a superset cannot increase the optimum. Hence
every completion $v$ of $u$ satisfies $LB(u)\leq\OPT(v)$.

The incumbent is the shortest feasible path found so far, of length $UB$.
The benchmark implementation starts from a discrete approximation and refines
its selected pieces with the convex solver. The compact production solver
starts with the first piece of each polygon and falls back to dynamic
programming over sampled polygon vertices if necessary.

\section{Pseudocode}

Algorithm~\ref{alg:bnb} mirrors the implemented search. The queue policy may
be FIFO or priority-based without affecting correctness. The compact
production implementation uses FIFO; the instrumented benchmark explores
nodes depth first and orders children using the incumbent path. Node order
affects performance, not the validity of pruning.

\begin{algorithm}[H]
\caption{Branch and bound for fixed-order nonconvex TPP}
\label{alg:bnb}
\begin{algorithmic}[1]
\Require $s,t$; ordered polygons $P_1,\ldots,P_k$; call limit $M$
\Ensure best path found, and whether its optimality was proved
\For{$i\gets1$ to $k$}
	\State decompose $P_i$ into convex pieces $C_{i,1},\ldots,C_{i,m_i}$
	\State $H_i\gets\conv(\bigcup_q C_{i,q})$
\EndFor
\State $(\pi^*,UB)\gets\Call{InitialFeasiblePath}{s,t,C}$
\State $Q\gets\{()\}$; $calls\gets0$
\While{$Q\neq\varnothing$ and $calls<M$}
	\State remove $u=(j_1,\ldots,j_d)$ from $Q$
	\If{$d=k$}
		\State $\pi\gets\Call{ConvexTPPPath}{s,t,C_{1,j_1},\ldots,C_{k,j_k}}$
		\State $calls\gets calls+1$
		\If{$\operatorname{length}(\pi)<UB$}
			\State $(\pi^*,UB)\gets(\pi,\operatorname{length}(\pi))$
		\EndIf
	\Else
		\For{$q\gets1$ to $m_{d+1}$}
			\State $v\gets(j_1,\ldots,j_d,q)$
			\State $R\gets(C_{1,j_1},\ldots,C_{d,j_d},C_{d+1,q},H_{d+2},\ldots,H_k)$
			\State $LB\gets\Call{ConvexTPPLength}{s,t,R}$
			\State $calls\gets calls+1$
			\If{$LB\leq UB$}
				\State insert $v$ into $Q$
			\EndIf
		\EndFor
	\EndIf
\EndWhile
\State $exact\gets(Q=\varnothing)$
\State \Return $(\pi^*,exact,calls)$
\end{algorithmic}
\end{algorithm}

The implementation keeps ties ($LB=UB$). This is conservative and allows it
to exhaust all nodes that may contain another optimum. If only one optimal
value is needed, pruning on $LB\geq UB$ is valid once feasibility and numerical
tolerances are handled carefully.

\section{Why the algorithm is correct}

\paragraph{Coverage.}
Each feasible path visits at least one convex piece of each polygon. Therefore
at least one complete piece vector represents every feasible path. Unless an
ancestor is pruned, that vector is reached as a leaf.

\paragraph{Safe pruning.}
For every node $u$, the convex hulls used by $R(u)$ contain all pieces allowed
in its descendants. Therefore $LB(u)$ is no larger than the value of any
completion. If $LB(u)>UB$, no descendant can improve the incumbent.

\paragraph{Optimality on exhaustion.}
If the queue becomes empty, every complete choice was either solved or has an
ancestor safely pruned. The incumbent is thus globally optimal for the given
polygon order. If the call or time limit is reached first, the path remains
feasible but optimality has not been proved; the solver reports
\texttt{exact=false}.

\section{Experimental evaluation}
\label{sec:experiments}

\subsection{Corpus and protocol}

We used the 558 instances in \texttt{instances\_socg\_simplified.zip}, released
with the TSPN solver of Fekete, Kniep, Krupke, and Perk from TU Braunschweig.\footnote{Project repository:
\url{https://github.com/tubs-alg/TSPN-SoCG-2026}.}
The source instances describe unordered closed TSPN tours. Our converter keeps
the polygon order stored in each JSON file and turns the problem into a path:
$s$ and $t$ are respectively the lower-left and upper-right corners of the
instance bounding box. Thus, these experiments test our fixed-order path
solver on their geometry; they are \emph{not} a runtime comparison with the
German solver, which also optimizes the cyclic visit order.

The code was compiled in release mode and run on a MacBook Pro with an Apple
M4 Pro. We used 12 worker threads, one repetition, no polygon or branching
limit, a 120-second guard per instance, and at most $10^6$ convex-solver calls per
instance. Optional experimental bounds were disabled, so all reported pruning
came from the convex-hull bound described above.

Table~\ref{tab:results} summarizes the complete run. A run is ``solved'' only
when the search tree was exhausted, which proves optimality for the converted
fixed-order instance.

\begin{table}[ht]
\centering
\caption{Results on the 558 converted TU Braunschweig instances.}
\label{tab:results}
\begin{tabular}{lr}
\toprule
Metric & Result \\
\midrule
Instances & 558 \\
Optimality proved & 517 (92.65\%) \\
Stopped at $10^6$ calls & 25 (4.48\%) \\
Stopped by the 120 s guard & 16 (2.87\%) \\
Total convex-solver calls & 39,246,477 \\
Bound calls & 37,744,960 (96.17\%) \\
Leaf calls & 1,374,884 (3.50\%) \\
Visited nodes & 13,815,892 \\
Pruned children & 23,928,516 (63.40\%) \\
Median calls per instance & 67 \\
90th percentile of calls & 114,791 \\
Median number of polygons & 30 \\
Maximum number of polygons & 60 \\
Median decomposed pieces & 50 \\
Maximum decomposed pieces & 503 \\
Median $\log_2$ of complete combinations & 16.49 \\
Maximum $\log_2$ of complete combinations & 138.72 \\
Convex solver share of measured work & 97.59\% \\
Wall-clock time with 12 workers & 727.34 s \\
\bottomrule
\end{tabular}
\end{table}

\subsection{Interpretation}

Three conclusions are immediate. First, the baseline proves optimality for
more than $92\%$ of this corpus, even though the largest formal search space
has approximately $2^{138.72}$ complete piece choices. The hull bound therefore
removes an enormous portion of the enumeration tree.

Second, $96.17\%$ of all convex calls compute bounds, while only $3.50\%$
solve leaves. Moreover, the convex solver accounts for $97.59\%$ of the
measured work. The principal bottleneck is therefore the number and cost of
bound solves, not convex decomposition or construction of the initial path.

Third, the median initial-to-final improvement is only $0.000168\%$, and the
median incumbent-to-final improvement is zero. The algorithm generally finds
a nearly optimal path early and spends most of its effort proving that no
better piece selection exists. This suggests prioritizing stronger or cheaper
bounds over a substantially more expensive initial heuristic.

The time guard is checked between convex calls, not inside the convex solver.
It is consequently not a hard wall-clock timeout: one instance spent 616.21 s
inside a single call before the B\&B could observe that the 120 s guard had
expired. This explains why the complete run took 727.34 s despite the nominal
per-instance limit and exposes a robustness issue to address in future runs.

These data also show why a single average is misleading: some geometrically
easy instances have enormous formal combination spaces but are eliminated near
the root, while a few instances require the full call budget. Both the
distribution and the worst cases must therefore be retained in future
comparisons.

\section{Current limitations and next experiments}

The present method is exact on exhaustion, but its worst-case search tree still
has $\prod_i m_i$ leaves. Its performance depends on four interacting factors:
the number of decomposition pieces, the quality of the initial incumbent, the
tightness of the hull relaxation, and the cost of each convex solve.

The instrumented benchmark already contains prototypes for piece-graph bounds,
port bounds, one-step refinement, contact-aware refinement, and safe piece
grouping. They were deliberately disabled in Section~\ref{sec:experiments} to
measure the baseline algorithm. The next controlled experiments should be:
\begin{enumerate}
	\item replace the fixed search order by best-bound search;
	\item improve the initial feasible solution with several cheap piece choices;
	\item use inexpensive combinatorial bounds before calling the convex solver;
	\item refine a hull only when the relaxed contact lies outside the real polygon;
	\item compare each addition separately by solved rate, convex calls, pruning
	rate, and wall-clock time.
\end{enumerate}

Extending the method to choose the polygon order is a separate combinatorial
layer. Noncrossing properties of optimal Euclidean tours may then provide
sequence-pruning rules, but they are not needed for the fixed-order correctness
result established here.

\section{Conclusion}

The algorithm reduces fixed-order nonconvex TPP to a discrete search whose
oracle is the exact convex TPP solver. Convex decomposition provides the
branches; convex hulls provide certified lower bounds; and any feasible path
provides an upper bound. This separation makes the proof short and the solver
auditable. The experimental question is no longer whether branch and bound can
avoid enumeration, but how to reduce the number and cost of the convex calls
needed to prove optimality on its hardest instances.

\end{document}
